\subsection*{Appendix-A.11: exact entrance derivative and a conditional axial-stress certificate}
The original interpolation (A.10), with \(a=1-\lambda\),
\(J_0=\log(1+\eta^2)\), \(\theta=1-\sigma(t/T_f)\) and
\(G=\exp((1-\theta)(J_0-\log2))\), implies
\[
r_f=\frac{e^{-aT_f}}{G(T_f,\eta)}
 \left[r_p+\int_0^{T_f}e^{at}G(t,\eta)\,dt\right].
\]
All pre-interpolation azimuthal stages retain the common
\((1+\eta^2)^{-1}\) factor. Thus \(\partial_\eta r_p=0\) and
differentiation under the integral yields exactly
\[
\partial_\eta r_f=-J_0'(\eta)\left[
\frac{2r_pe^{-aT_f}}{1+\eta^2}
+\int_0^{T_f}e^{-au}\left(\frac2{1+\eta^2}\right)^{\sigma(u/T_f)}
                  \sigma(u/T_f)\,du\right].
\]
Consequently
\[
|\partial_\eta r_f|\le 2r_pe^{-aT_f}
  +\frac2a[\sigma(1/4)+e^{-aT_f/4}].
\]
For \(T_f\ge512\), \(0<\lambda\le1/20\),
\(r_p\le1/a\), the right side is \(<1.41\cdot10^{-6}\).
The source step has \(\|\sigma'\|_\infty\le64\), so the
interpolation slope belongs to \([-\lambda-1/10,-\lambda]\).

Let the prescribed normalized \(\sigma'\) bumps be centered
at \(L-d_i\), \(L=30\log(1/\lambda)\),
\(d_1=3,d_2=1\), and let \(E=E_0(1+q)\),
\(q=\sum_ic_i(\eta)\beta_i\). Put \(b=1+2\lambda\),
\[
u_i=e^{-ad_i}\int e^{az}\beta(z)\,dz,\quad
v_i=e^{bd_i},\quad
\kappa=\frac{\int e^{-bz}\beta(z)^2\,dz}
             {2\int e^{-bz}\beta(z)\,dz},\quad
\epsilon=(1/a-r_f)e^{-aL}.
\]
The TWO actual A.11 requirements are precisely
\[
u_1c_1+u_2c_2=\epsilon,\qquad
v_1c_1+v_2c_2+\kappa(v_1c_1^2+v_2c_2^2)=0.
\]
The support and exponential-weight estimates give
\[
|\det B|>5,\quad \|B^{-1}\|_\infty<7,\quad
\kappa<150,\quad v_1+v_2<36.
\]
Source comparison gives \(|\epsilon|\le3\lambda^{30a}\),
\(|\partial_\eta\epsilon|\le5\lambda^{30a}\).
On the ball \(\|c\|_\infty\le14|\epsilon|\), the exact
quadratic moment map has Lipschitz constant at most
\(7\cdot10800\cdot14|\epsilon|<1/2\). It has a unique
small solution with
\[
\|c\|_\infty\le42\lambda^{30a},\qquad
\|\partial_\eta c\|_\infty\le70\lambda^{30a}.
\]
The source \(\sigma'\)-bump satisfies \(\|\beta\|_\infty<214\),
\(\|\beta'\|_\infty<56000\); hence
\(\|q\|_\infty<1/4\), \(\|q_\eta\|_\infty<1\) and
\(\|q_y\|_\infty<3\lambda/8\).
Thus the corrected log slope is at most \(-\lambda/2\).

The entrance-moment derivative alone DOES NOT bound \(N_s/E\).
Using the exact source identity (4.16), and explicitly assuming
the completed amplitude constraint \(S(\infty)=0\) and
the pressure gauge \(\Pi(\infty)=0\), one obtains after the
pulse (\(U=M=J=0\)):
\[
N_s(y,\eta)=2\eta\int_y^\infty
(h e^{v-y}-A)E(v,\eta)^2\,dv
+\frac{1-\eta^2}{2}\int_y^\infty
(1-e^{v-y})\partial_\eta E(v,\eta)^2\,dv .
\]
If the remaining source slopes obey \(\ell_0\le0\),
the \(-1\) hold lasts \(4\log(1/h)\), and
\(\ell_0\le-3h/4\) afterward, then the entire future
baseline energy obeys
\[
\int_y^\infty E_0(v)^2\,dv\le E_0(y)^2,\quad
\int_y^\infty e^{v-y}E_0(v)^2\,dv
\le(93/20+2h^7/3)E_0(y)^2
\]
on both A.11 supports. Since their union has length \(3/5\),
the exact \(N_s\) formula and bounds on \(q,q_\eta\) give
\[
\frac{|N_s|}{E}\le4E_0
\le40P_*\lambda^{15}e^{-13/(2\lambda)}.
\]
Thus the sought \(\sup|N_s|/E\le\lambda/2\) follows from
the explicit sufficient scalar gate
\[
80P_*\lambda^{14}e^{-13/(2\lambda)}\le1.
\]
This is a CONDITIONAL SOURCE-FAMILY statement. The source
requires \(S(\infty)=0\), actual post-pulse slope bounds, and
one simultaneously admissible parameter realization; the
diagnostic numbers in the accompanying exact checker do not
certify those missing hypotheses. No unforced Theorem 1.1,
global singularity, or end-to-end Lean proof is asserted.
